\DocumentMetadata{
  pdfversion=2.0,
  pdfstandard=ua-2,
  lang=en-US,
  tagging=on,
  tagging-setup={math/setup=mathml-SE,tabsorder=structure}
}
\documentclass[11pt,letterpaper]{article}
\usepackage[margin=0.68in,includefoot]{geometry}
\usepackage{newcomputermodern}
\usepackage{amsmath,amscd,mathtools}
\usepackage{tikz,adjustbox}
\providecommand{\square}{\mdlgwhtsquare}
\usepackage[hidelinks]{hyperref}
\hypersetup{
  pdftitle={Analysis PhD Exam, August 2007},
  pdfauthor={University of Florida Department of Mathematics},
  pdfsubject={Graduate mathematics examination},
  pdfdisplaydoctitle=true,
  bookmarksopen=true
}
\setcounter{secnumdepth}{0}
\setlength{\parindent}{0pt}
\setlength{\parskip}{0.28em}
\setlength{\emergencystretch}{3em}
\setlength{\leftmargini}{2.15em}
\setlength{\leftmarginii}{2.65em}
\setlength{\leftmarginiii}{3em}
\pagestyle{empty}
\raggedbottom
\allowdisplaybreaks
\NewTaggingSocketPlug{block/list/label}{examproblem}{%
  \tagstructbegin{tag=Lbl}\tagstructend%
  \tagstructbegin{tag=\LBody}%
  \tagstructbegin{tag=\ProblemHeadingTag,title-o={\ProblemHeadingText}}%
  \tagmcbegin{tag=\ProblemHeadingTag,actualtext={\ProblemHeadingText}}%
  #2%
  \tagmcend%
  \tagstructend%
}
\newenvironment{examproblems}[2]{%
  \def\ProblemHeadingTag{#2}%
  \AssignTaggingSocketPlug{block/list/label}{examproblem}%
  \begin{enumerate}%
  \setcounter{enumi}{#1}%
  \renewcommand{\labelenumi}{\textbf{\theenumi.}}%
  \setlength{\topsep}{0.35em}%
  \setlength{\partopsep}{0pt}%
  \setlength{\itemsep}{0.48em}%
  \setlength{\parsep}{0.18em}%
}{\end{enumerate}}
\newenvironment{examsubparts}{%
  \AssignTaggingSocketPlug{block/list/label}{default}%
  \begin{enumerate}%
  \setlength{\topsep}{0.18em}%
  \setlength{\partopsep}{0pt}%
  \setlength{\itemsep}{0.16em}%
  \setlength{\parsep}{0.1em}%
}{\end{enumerate}}
\newenvironment{examinstructions}{%
  \par\begingroup\itshape
}{%
  \par\endgroup\vspace{0.18em}
}
\makeatletter
\renewcommand\subsection{\@startsection{subsection}{2}{0pt}%
  {-1.25ex plus -.25ex minus -.1ex}%
  {0.55ex plus .1ex}%
  {\normalfont\large\bfseries}}
\renewcommand{\@maketitle}{%
  \newpage\null\vskip -1em%
  {\footnotesize\bfseries
    \href{https://www.ufl.edu/}{University of Florida}, \href{https://math.ufl.edu/}{Department of Mathematics}\par}%
  \vskip -0.5em%
  \tagstructbegin{tag=H1,title={Analysis PhD Exam, August 2007}}%
  \tagpdfparaOff%
  \tagmcbegin{tag=H1}%
    {\LARGE\bfseries\href{https://gma.math.ufl.edu/exams/analysis-phd/}{Analysis PhD Exam}, August 2007\par}%
  \tagmcend%
  \tagpdfparaOn%
  \tagstructend%
  \vskip 0.25em%
  \tagmcbegin{artifact}%
    \hrule height 1pt%
  \tagmcend%
  \vskip 1em%
}
\makeatother
\title{Analysis PhD Exam, August 2007}
\author{}
\date{}
\begin{document}
\maketitle
\thispagestyle{empty}
\begin{examinstructions}
Attempt EIGHT questions. To receive credit be sure to state clearly the theorems used in solutions.
\end{examinstructions}
\begin{examproblems}{0}{H2}
\def\ProblemHeadingText{Problem 1.}
\item
Let \((X,\mathcal S,\mu)\) be the following measure space: \(X\) is the set of positive integers, \(\mathcal S\) is the power set of~\(X\), and~\(\mu\) is the counting measure. If~\(f\) is~\(\mu\)-integrable, does \(\sum |f(m)|\) converge? Prove. Is the converse true? Prove.
\def\ProblemHeadingText{Problem 2.}
\item
Let~\(\mu\) and~\(\nu\) be finite valued signed measures such that \(\mu\ll\nu\) and \(\nu\ll\mu\). Prove that \(\left(\frac{d\mu}{d\nu}\right)\left(\frac{d\nu}{d\mu}\right)=1\) a.e.
\def\ProblemHeadingText{Problem 3.}
\item
Let~\(\Sigma\) be an algebra of subsets of a set~\(S\) and \(\sigma(\Sigma)\) the σ-algebra generated by~\(\Sigma\). State the Hahn extension theorem for a σ-finite measure~\(\mu\) defined on~\(\Sigma\). Use a Monotone Class argument to prove that the extension is unique.
\def\ProblemHeadingText{Problem 4.}
\item
Let \(\{x_m\}\) be a sequence of elements from a Banach space \(\mathcal X\). Suppose \(\sum \|x_m\|<\infty\). Does \(\sum x_m\) converge unconditionally? Prove.
\def\ProblemHeadingText{Problem 5.}
\item
Let \(\mathcal M\) and \(\mathcal N\) be closed linear subspaces of a Hilbert space \(\mathcal H\). Assume \(\mathcal M\perp\mathcal N\). Show \(\mathcal M+\mathcal N\) is closed.
\def\ProblemHeadingText{Problem 6.}
\item
Let \(\mathcal X\) be a Banach space and suppose \(\Pi:\mathcal X\to\mathcal X^{**}\) is the natural embedding map. Is \(\Pi(\mathcal X)\) closed in \(\mathcal X^{**}\)? Prove.
\def\ProblemHeadingText{Problem 7.}
\item
Let \(L^1(\mathbb R)\) be the space of Lebesgue integrable functions on \(\mathbb R\). Assume~\(E\) is a Lebesgue measurable set. Denote by~\(Y\) the set of all functions \(f\in L^1(\mathbb R)\) such that \(f(t)=f(-t)\) for almost all \(t\in E\). Prove~\(Y\) is closed in \(L^1(\mathbb R)\).
\def\ProblemHeadingText{Problem 8.}
\item
Let \((\Omega,\mathcal F,P)\) be a probability space and \(\mathcal D\) a sub-σ-algebra of \(\mathcal F\). Suppose \(f\in L^2(\Omega,\mathcal F,P)\). Show by Hilbert space theory that there exist a \(\mathcal D\)-measurable function~\(h\) such that
\[
\int_D f\,dP=\int_D h\,dP,\qquad D\in\mathcal D.
\]
\def\ProblemHeadingText{Problem 9.}
\item
Let \(\mathcal Y\) and \(\mathcal Z\) be closed linear subspaces of the Banach space \(\mathcal X\) such that each \(x\in\mathcal X\) has a unique representation \(x=y+z\), where \(y\in\mathcal Y\) and \(z\in\mathcal Z\). Show that there exists a constant~\(K\) such that \(\|y\|\le K\|x\|\) and \(\|z\|\le K\|x\|\), for all \(x\in\mathcal X\), where \(x=y+z\) is the above representation of~\(x\).
\end{examproblems}
\end{document}
